% Chapter 2

\chapter{Estado da Arte} % Main chapter title
\label{chap:Chapter2} % For referencing the chapter elsewhere, use \ref{chap:Chapter2}

% NOTA: capítulo em desenvolvimento. Não inserir resultados, contagens ou lacunas
% antes de concluir a pesquisa, a seleção e a síntese dos estudos.
%
% A CONFIRMAR ANTES DE ENTREGAR (procura por "TODO"):
%  - As RQ-R1 e RQ-R2 abaixo têm de ser iguais às RQ1 e RQ2 do Capítulo 1.
%  - Fonte Scopus: ainda não foi pesquisada (decidir Scopus vs. ScienceDirect).
%  - Q3 (segurança) ainda não foi testada no piloto.
%  - Número de revisores na triagem.
%  - Cada [CITAR] precisa de uma referência verificada.

%----------------------------------------------------------------------------------------
\section{Introdução}
\label{sec:chap2_intro}

Este capítulo apresenta uma revisão sistemática da literatura sobre sistemas conversacionais baseados em \glspl{LLM} cujas respostas dependem de um contexto definido pela aplicação. A revisão organiza-se em três dimensões: as abordagens técnicas para fornecer esse contexto ao modelo e restringir as respostas; os métodos de avaliação da fidelidade ao contexto, da correção factual e da recusa de perguntas fora do âmbito; e os compromissos operacionais reportados, nomeadamente fiabilidade, segurança, latência, custo e integração.

No caso de estudo desta dissertação, o contexto associado à pergunta em curso no Adhara QBank é fornecido diretamente ao modelo, sem uma etapa de recuperação de informação. A literatura sobre \emph{retrieval-augmented generation} (RAG) é, por isso, considerada como termo de comparação e como principal fonte de métodos de avaliação, e não como solução assumida. A síntese será relacionada com esse caso de estudo: um assistente para esclarecer dúvidas sobre a pergunta em curso.

O capítulo apresenta o protocolo da revisão (Secção~\ref{sec:chap2_protocol}), os resultados da pesquisa e da seleção (Secção~\ref{sec:chap2_results}), a síntese crítica dos estudos (Secção~\ref{sec:chap2_synthesis}) e a lacuna de investigação que daí resulta (Secção~\ref{sec:chap2_gap}).

%----------------------------------------------------------------------------------------
\section{Protocolo da Revisão}
\label{sec:chap2_protocol}

A revisão segue o PRISMA 2020 \textbf{[CITAR]}. O protocolo foi definido antes da pesquisa final: questões, fontes, expressões de pesquisa, critérios de seleção, avaliação da qualidade e campos de extração. As alterações feitas depois da pesquisa piloto estão justificadas na Secção~\ref{sec:chap2_search}.

\subsection{Questões da Revisão}
\label{sec:chap2_rqs}

\begin{description}
    \item[RQ-R1] Que abordagens utilizam os sistemas conversacionais baseados em \glspl{LLM} para limitar as respostas a um contexto específico da aplicação?
    \item[RQ-R2] Como são avaliadas, nestes sistemas, a fidelidade ao contexto fornecido, a correção factual e a recusa de perguntas fora do âmbito?
\end{description}

% TODO: copiar estas RQ para o Capítulo 1 (secção Questões de Investigação).

Os compromissos operacionais (latência, custo, segurança e integração) não constituem uma questão autónoma: são extraídos de cada estudo, quando reportados, e discutidos na síntese (Secção~\ref{sec:chap2_synthesis}).

\subsection{Fontes e Estratégia de Pesquisa}
\label{sec:chap2_search}

As fontes são Scopus, IEEE Xplore e ACM Digital Library. O Scopus é uma base multidisciplinar que indexa vários editores; o IEEE Xplore e a ACM Digital Library cobrem a literatura de computação e engenharia de software. Foram escolhidas antes da pesquisa final. O \emph{snowballing}, se utilizado, será documentado como método complementar e não misturado com os resultados das bases. % TODO: confirmar Scopus (ver nota no topo).

\paragraph{Pesquisa piloto.}
Antes de fixar as expressões, foi feita uma pesquisa piloto em IEEE Xplore, ACM Digital Library e ScienceDirect, limitada aos primeiros 30 títulos de cada pesquisa e sem acesso aos resumos. O piloto serviu apenas para ajustar as expressões e não entra nas contagens finais. Levou a três decisões: (i) estreitar as expressões a sistemas conversacionais e de resposta a perguntas, porque as versões gerais devolviam sobretudo \emph{hardware} de inferência, geração de código e recomendação; (ii) retirar latência e custo dos blocos de pesquisa e tratá-los como campos de extração, pelo mesmo motivo; (iii) tratar as revisões e os \emph{surveys}, muito frequentes nos resultados, pelo critério de exclusão E4.

\paragraph{Expressões de pesquisa.}
Cada expressão combina blocos de conceitos: os sinónimos de um bloco são ligados por \texttt{OR} e os blocos por \texttt{AND}. Não utilizam truncagem com \texttt{*} e respeitam o limite de conectores booleanos por campo. Pesquisam-se o título, o resumo e as palavras-chave, na sintaxe de cada fonte (no Scopus, \texttt{TITLE-ABS-KEY(...)}). A versão executada em cada fonte é guardada literalmente.

\paragraph{Q1 --- contexto e sistemas conversacionais (RQ-R1).}
\begin{verbatim}
("large language model" OR LLM)
AND (grounding OR "context injection" OR "prompt context"
     OR "retrieval-augmented generation")
AND (chatbot OR "conversational system" OR "question answering")
\end{verbatim}

\paragraph{Q2 --- fidelidade, correção e recusa fora do âmbito (RQ-R2).}
\begin{verbatim}
("large language model" OR LLM)
AND (faithfulness OR factuality OR "out-of-scope")
AND (chatbot OR "conversational system")
AND (evaluation OR benchmark)
\end{verbatim}

\paragraph{Q3 --- segurança do âmbito (RQ-R1).}
\begin{verbatim}
("large language model" OR LLM)
AND ("prompt injection" OR "prompt leakage" OR jailbreak)
AND (chatbot OR "conversational system")
\end{verbatim}
% TODO: Q3 ainda não foi testada no piloto; ajustar ou retirar conforme o resultado.

\begin{table}
\caption{Pesquisas e questões da revisão.}
\label{tab:queries}
\centering
\small
\begin{tabular}{p{1.4cm} p{3cm} p{8.6cm}}
\toprule
\tabhead{Pesquisa} & \tabhead{Questão} & \tabhead{Objetivo} \\
\midrule
Q1 & RQ-R1 & Abordagens que fornecem contexto ao modelo e limitam as respostas \\
Q2 & RQ-R2 & Avaliação de fidelidade, correção e recusa fora do âmbito \\
Q3 & RQ-R1 & Mecanismos de segurança que protegem o âmbito (injeção e fuga de instruções) \\
\bottomrule
\end{tabular}
\end{table}

\paragraph{Regra de paragem e registo.}
Em cada fonte e para cada pesquisa são analisados os primeiros 30 resultados, ordenados por relevância; este limite é a regra de paragem. Para cada par (fonte, pesquisa) regista-se a expressão exata, a data de execução, o intervalo temporal, a língua, o total apresentado pela fonte e o número de registos analisados, numa tabela de registo das pesquisas. Os registos são exportados com resumo e importados no Rayyan \textbf{[CITAR]}.

\subsection{Critérios de Inclusão e Exclusão}
\label{sec:chap2_criteria}

\begin{table}
\caption{Critérios de inclusão e exclusão da revisão.}
\label{tab:criteria}
\centering
\small
\begin{tabular}{p{1cm} p{5.8cm} p{1cm} p{5.8cm}}
\toprule
\tabhead{} & \tabhead{Inclusão} & \tabhead{} & \tabhead{Exclusão} \\
\midrule
I1 & Publicação entre 27-10-2021 e 27-10-2026, inclusive & E1 & Duplicados \\
I2 & Língua inglesa ou portuguesa & E2 & Editoriais, opiniões, livros e registos sem informação metodológica suficiente \\
I3 & Estudo primário sobre sistema conversacional ou de resposta a perguntas baseado em \gls{LLM}, cujas respostas dependem de um contexto definido (RAG, injeção no \emph{prompt} ou base de conhecimento) & E3 & Sem \gls{LLM}, sem sistema conversacional ou de resposta a perguntas, ou sem relação com respostas contextualizadas nem com a sua avaliação \\
I4 & Método, sistema ou avaliação descritos com detalhe suficiente para extrair dados & E4 & Revisões e \emph{surveys}, excluídos da síntese principal; podem apoiar o enquadramento e o \emph{snowballing} \\
\bottomrule
\end{tabular}
\end{table}

O critério I3 não exige recuperação de informação: um estudo em que o contexto é fornecido diretamente ao modelo cumpre-o tal como um estudo com RAG.

\subsection{Seleção dos Estudos}
\label{sec:chap2_selection}

A seleção tem três etapas: identificação e remoção de duplicados; triagem de títulos e resumos; e avaliação dos textos integrais. A triagem é feita no Rayyan \textbf{[CITAR]}, com as decisões e os motivos de exclusão registados. Na avaliação dos textos integrais, cada exclusão é associada a um dos critérios E1 a E4. O processo é documentado num diagrama PRISMA \textbf{[CITAR]}, que indica os registos por fonte, os duplicados removidos, os excluídos em cada etapa e os estudos incluídos. % TODO: indicar o número de revisores (por exemplo, segunda triagem pelo orientador).

\subsection{Avaliação da Qualidade}
\label{sec:chap2_quality}

A elegibilidade (critérios de seleção) é distinta da qualidade. Os estudos incluídos são avaliados com cinco critérios, cada um classificado como \emph{sim} (1), \emph{parcial} (0,5) ou \emph{não} (0):
\begin{enumerate}
    \item[QA1] O método ou sistema está descrito com clareza suficiente para ser compreendido.
    \item[QA2] O contexto fornecido ao modelo e a sua origem estão descritos.
    \item[QA3] A avaliação usa dados, métricas ou comparações explícitos.
    \item[QA4] Os resultados são apresentados de forma verificável (dados, dimensão da amostra, avaliadores).
    \item[QA5] São discutidas limitações ou ameaças à validade.
\end{enumerate}
A pontuação não exclui estudos: serve para ponderar a evidência na síntese.

\subsection{Extração de Dados}
\label{sec:chap2_extraction}

Para cada estudo incluído são extraídos os campos da Tabela~\ref{tab:extraction}. A matriz completa fica em apêndice; no texto apresenta-se apenas uma tabela de classificação.

\begin{table}
\caption{Campos de extração de dados.}
\label{tab:extraction}
\centering
\small
\begin{tabular}{p{0.6cm} p{3.4cm} p{7.2cm} p{1.6cm}}
\toprule
\tabhead{\#} & \tabhead{Campo} & \tabhead{Descrição} & \tabhead{Questão} \\
\midrule
1 & Referência e ano & Autores, ano e tipo de publicação & --- \\
2 & Domínio e aplicação & Educação, saúde, apoio ao cliente, entre outros & RQ-R1 \\
3 & Modelo e abordagem & Modelo de linguagem; \emph{prompting}, RAG, \emph{fine-tuning} ou combinação & RQ-R1 \\
4 & Chegada do contexto & Como o contexto chega ao modelo: RAG, injeção direta no \emph{prompt}, base de conhecimento ou outro & RQ-R1 \\
5 & Controlo e recusa & Mecanismos para limitar as respostas e recusar perguntas fora do âmbito & RQ-R1, RQ-R2 \\
6 & Avaliação & Método (automática, \emph{LLM-as-a-judge}, humana), dados, métricas e participantes & RQ-R2 \\
7 & Resultados e limitações & Principais resultados e limitações indicadas pelos autores & RQ-R2 \\
8 & Aspetos operacionais & Latência, custo, segurança e integração, quando reportados & RQ-R1 \\
\bottomrule
\end{tabular}
\end{table}

\subsection{Ameaças à Validade}
\label{sec:chap2_threats}

Esta revisão tem limitações conhecidas. O limite de 30 resultados por fonte e pesquisa, ordenados por relevância, pode deixar de fora estudos relevantes que a fonte classifica mais abaixo; por isso o total apresentado por cada fonte é registado. A restrição ao inglês e ao português e às três fontes escolhidas limita a cobertura. A literatura sobre \glspl{LLM} evolui depressa e a janela temporal termina em 27-10-2026. % TODO: acrescentar o número de revisores quando definido.

%----------------------------------------------------------------------------------------
\section{Resultados da Pesquisa}
\label{sec:chap2_results}

A pesquisa e a seleção dos estudos ainda estão em curso. Por esse motivo, este capítulo não apresenta contagens nem um diagrama de fluxo como resultados finais. Após a execução final das pesquisas, serão registados os resultados por fonte, a remoção de duplicados, as decisões de seleção por título e resumo, a avaliação do texto integral e os estudos incluídos. O diagrama PRISMA será inserido depois de estes valores serem confirmados.

% \begin{figure}[t]
% \centering
% \includegraphics[width=0.8\textwidth]{ch2/assets/prisma}
% \caption[Diagrama PRISMA]{Diagrama PRISMA do processo de seleção.}
% \label{fig:prisma}
% \end{figure}

% Tabela de classificação dos estudos (preencher só depois da extração):
% \begin{table}
% \caption{Classificação dos estudos incluídos.}
% \label{tab:classification}
% \centering
% \small
% \begin{tabular}{p{3cm} p{2.5cm} p{3cm} p{2.5cm} p{2cm}}
% \toprule
% \tabhead{Estudo} & \tabhead{Domínio} & \tabhead{Abordagem} & \tabhead{Contexto} & \tabhead{Avaliação} \\
% \midrule
% \bottomrule
% \end{tabular}
% \end{table}

% Tabela de registo das pesquisas (uma linha por fonte e pesquisa):
% fonte | pesquisa | expressão executada | data | filtros | total na fonte | analisados

%----------------------------------------------------------------------------------------
\section{Síntese Crítica}
\label{sec:chap2_synthesis}

A síntese deve agrupar, comparar, explicar e criticar os estudos, e não resumi-los um a um. A frequência de uma técnica descreve a sua utilização, não prova que seja a mais eficaz.

\subsection{Abordagens}
% TODO: agrupar por como o contexto chega ao modelo (RAG, injeção direta, base de conhecimento) e por técnica (prompting, fine-tuning); comparar. Verificar quantos estudos usam injeção direta, sem recuperação.

\subsection{Avaliação}
% TODO: como são avaliadas a fidelidade, a correção e a recusa; métodos, métricas e participantes; evidência de impacto na aprendizagem, se existir.

\subsection{Controlo de Âmbito e Segurança}
% TODO: injeção de instruções, fuga do prompt, recusa fora de âmbito.

\subsection{Compromissos Operacionais}
% TODO: latência, custo e integração reportados; como se relacionam com a abordagem usada.

\subsection{Limitações Reportadas}
% TODO: limitações recorrentes e contradições na evidência.

%----------------------------------------------------------------------------------------
\section{Lacuna de Investigação}
\label{sec:chap2_gap}

% TODO: escrever só depois da síntese, seguindo a cadeia
% Evidência -> Limitação -> Lacuna -> Questão de investigação -> Contribuição.
% Hipótese de trabalho a verificar (NÃO escrever como resultado): a literatura de respostas
% fundamentadas em contexto concentra-se em RAG; pode haver pouca evidência sobre contexto
% fechado por item, injetado diretamente, com avaliação de fidelidade e de recusa.

%----------------------------------------------------------------------------------------
\section{Conclusões do Capítulo}
\label{sec:chap2_conclusions}

% TODO: o que se sabe, o que falta saber e como o trabalho desta dissertação responde à lacuna.